\documentclass[10pt,conference,compsocconf]{IEEEtran}

\usepackage{hyperref}
\usepackage{graphicx}
\usepackage{booktabs}

\begin{document}
\title{How Unified is Your Encoder?\\ Measuring Representational Convergence in Any-to-Any Vision Models%
\thanks{Code and configs: \url{https://github.com/arthur-mrgt/one-trunk-or-many/tree/null-hypothesis}}}

\author{
  Arthur Margeat (330258), Adrien Clement (345535),
  Albert Fares (341018), Martina Gatti (341013)\\
  \textit{CS-503 Visual Intelligence --- Project Progress Report 1 (Revised)}
}

\maketitle

\section{Milestone Progress}

\subsection{Problem and milestone overview}
Any-to-any models such as 4M~\cite{4m} and its successor
4M-21~\cite{4m21} process every modality through a single shared
Transformer~\cite{transformer} encoder, with a separate learned input
embedding per modality. Whether this shared trunk produces a unified
representation, or merely routes inputs through effectively disjoint
modality pathways~\cite{zamir}, is open. The question matters more
broadly: it directly bears on hypotheses such as platonic
representational convergence~\cite{platonic} and on observations of
persistent modality gaps in multi-modal models~\cite{modality_gap}.

RQ1 of our proposal asks whether 4M's shared encoder maps different
modalities of the same scene into a common representation, and whether
the signal is consistent across three complementary similarity views.
\textbf{CKA}~\cite{cka} captures \emph{global geometric} structure: it
checks whether two layer representations preserve the same pairwise
similarities between scenes; it operates on Gram matrices, so it is
robust to the activation dimension. \textbf{PWCCA}~\cite{pwcca}
captures \emph{directional} alignment: built on canonical correlation
analysis, it tests whether the two representations align along the
same principal directions, after a low-dimensional projection.
\textbf{k-NN overlap} captures \emph{local topological} similarity:
``topological'' here means we only care about which scenes are close
to which, not about absolute distances or directions, so the metric
asks whether the nearest neighbours of a scene are the same in both
representations. Together these three views cover global geometry,
directional structure, and local neighbourhood preservation. We
measure each metric layer by layer and compare it against an empirical
null built from mismatched scenes, so that observed similarity is read
relative to a noise baseline rather than in absolute terms.

For milestone~1, we built a minimal proof-of-concept (POC) of this
protocol on one Hypersim~\cite{hypersim} subset and one modality pair
(RGB-Depth). The aim was end-to-end pipeline coherence before scaling
to all pairs and a larger sample budget.

\subsection{What the POC computed}
The POC ran four conceptual stages.
(i) For every scene, RGB and depth were independently forwarded
through the 4M-7B encoder, recording activations at each of the $12$
encoder blocks.
(ii) For every layer, CKA, PWCCA, and k-NN overlap were computed
between matched RGB and depth activations.
(iii) For PWCCA and k-NN, which need a low-dimensional input, a
single PCA was fit jointly on the stacked RGB and depth activations
and both modalities were projected through the same components.
(iv) A null distribution was built by sampling pairs of distinct
scenes $(A,B)$ and recomputing each metric between RGB of $A$ and
depth of $B$, using all samples in those scenes; about $1000$ such
mismatched draws were aggregated per layer, and per-layer p-values
were obtained by tail counts.

\subsection{POC results}
The POC ran end-to-end on $4325$ matched RGB-depth samples, $1000$
null draws, and the $12$ encoder blocks of 4M-7B-CC12M.
\textbf{CKA} rose roughly monotonically from $\approx 0.10$ at
layer~$0$ to $\approx 0.19$ at layer~$10$, before a small drop to
$\approx 0.17$ at the last block (Fig.~\ref{fig:metric-vs-layer-v2}).
\textbf{k-NN overlap} rose from $\approx 0.09$ to $\approx 0.35$
around layer~$9$. \textbf{PWCCA} disagreed: matched values stayed in
$0.25$--$0.56$ while the null mean was $\approx 0.70$ at every layer,
so all PWCCA p-values were above $0.8$
(Fig.~\ref{fig:pvalue-vs-layer-v2}).

\begin{figure}[t]
  \centering
  \includegraphics[width=\columnwidth]{figures/metric_vs_layer.pdf}
  \caption{\small POC layer-wise similarity (RGB-Depth, $n{=}4325$).}
  \label{fig:metric-vs-layer-v2}
\end{figure}

\begin{figure}[t]
  \centering
  \includegraphics[width=0.65\columnwidth]{figures/pvalue_vs_layer.pdf}
  \caption{\small POC per-layer p-values, shown only as diagnostics:
  the underlying null is biased (Sec.~\ref{sec:errors}).}
  \label{fig:pvalue-vs-layer-v2}
\end{figure}

\subsection{Why these results are not directly interpretable}
\label{sec:errors}
The CKA and k-NN trends move in the direction predicted by our
proposal, which is encouraging as a pipeline sanity check. However,
the POC kept two methodological simplifications that affect the
significance assessment.

\paragraph{(E1) Null built from whole scene pairs}
Each null draw used two whole mismatched scenes and computed the
metric on all samples available in those scenes. As a result, the
matched score and each null score are not estimated on comparable
sample sizes, and intra-scene structure leaks into the null. The
correct construction is to draw a fixed-size set of mismatched
single-image pairs whose size matches the matched estimate.

\paragraph{(E2) Joint PCA across the two modalities}
For PWCCA and k-NN, a single PCA was fit on both modalities at once.
This finds directions of joint variance that can artificially align
the two spaces, biasing CCA-based PWCCA in particular. The standard
practice is to reduce each modality independently before comparing.

Together, (E1) and (E2) explain the PWCCA anomaly (matched values
below the null mean) and make the absolute p-values for CKA and k-NN
preliminary as well, even where the trend itself looks plausible.

\subsection{Current state: full RQ1 pipeline ready}
After the POC, we corrected the two simplifications and finalised the
RQ1 pipeline:
\begin{itemize}\itemsep0pt
  \item null built from image-level mismatched pairs with per-draw
        size matched to the matched estimate (fixes E1);
  \item PCA fit independently per modality for PWCCA and k-NN
        (fixes E2);
  \item adaptive null sampling with multiple-testing correction via
        BH-FDR~\cite{bh1995} and intermediate checkpointing, so longer
        runs are robust to interruption;
  \item configuration-driven support for all proposal modality pairs
        and for both Hypersim and DIODE~\cite{diode}, with the null
        computed for each metric.
\end{itemize}
The pipeline is ready to run the full RQ1 protocol; we did not yet
execute it at full scale due to time and GPU-availability constraints
before submission.

\subsection{RQ1 blindspots and motivation for RQ2}\label{sec:rq2-motivation}
Even when the corrected RQ1 pipeline is run end-to-end, the protocol
only answers a specific question: are the per-layer representations of
different modalities geometrically and topologically similar? This
leaves two blindspots that RQ2 is designed to complement.

(i)~\emph{Functional compatibility.} Two representations can share
global geometry while still encoding information in incompatible
formats. A high CKA, for instance, only says that pairwise distances
between scenes are preserved across modalities; it does not say that
the encoder uses the same internal code on both sides. In other words,
RQ1 tells us whether the spaces have the same shape, not whether they
``speak the same language''. (ii)~\emph{Ceiling ambiguity.} Even when
significant convergence is detected, RQ1 cannot decide whether the
variance left unaligned is legitimate modality-specific information
(e.g.\ colour for RGB, absolute scale for depth) or hidden shared
information that our metrics simply fail to capture. The metrics
therefore give a lower bound on unification, not an upper one.

RQ2 is included as an exploratory complement, not a fixed protocol.
Its goal is to probe these two questions beyond pure similarity
measurements (e.g.\ via cross-modal transfer or residual analyses
inspired by~\cite{4m,4m21}), reusing the activations and null
artefacts already produced by RQ1. The exact protocols will be chosen
once RQ1 results are stable, so that RQ2 effort is directed where the
RQ1 outcome makes it most informative.

\section{Discussion \& Next Steps}\label{sec:next}

\paragraph{Reading}
The implementation deliverable is complete: a reusable RQ1 benchmark
with a corrected null and reduction protocol, plus a clear path from
RQ1 outputs to RQ2 protocols. The POC indicates consistent trends for
CKA and k-NN, while final significance is intentionally deferred until
the corrected pipeline is run end-to-end.

\paragraph{Risks}
(i)~\emph{Confounders}: similarity may partly come from sources
upstream of the shared encoder, in particular per-modality input
embeddings~\cite{4m} and shared positional encodings~\cite{transformer},
rather than from encoder unification itself.
(ii)~\emph{Dataset specificity}: Hypersim is synthetic; results need to
be checked on real data (DIODE~\cite{diode}).
(iii)~\emph{Scale}: null behaviour and significance can shift with
larger sample budgets, so milestone-1 numbers are not yet
load-bearing.

\paragraph{Action items before the next report}
\begin{enumerate}\itemsep0pt
  \item Run the corrected RQ1 pipeline at full scale on Hypersim,
        across all targeted modality pairs, with the adaptive null
        per metric.
  \item Replicate on DIODE~\cite{diode} for synthetic-to-real
        robustness.
  \item Quantify the main confounders identified above, isolating the
        per-modality embedding effects upstream of the shared encoder.
  \item Once RQ1 results are stable, start the RQ2 exploration
        outlined in Sec.~\ref{sec:rq2-motivation}, choosing the
        concrete protocols based on what RQ1 reveals.
\end{enumerate}

\section{Author Contribution Statement}
\textit{Arthur} led the extraction stage and the 4M model wrapper,
implemented the corrected final RQ1 pipeline, and contributed to the
write-up. \textit{Albert} and \textit{Martina} implemented the three
metrics (CKA, PWCCA, k-NN overlap) and the null-hypothesis stage.
\textit{Adrien} contributed to results analysis and the write-up. All
authors participated in shaping the research question.

\bibliographystyle{IEEEtran}
\bibliography{literature}

\end{document}
